\centerline{\Huge Условия и решения задач.}
\centerline{\bf \Huge 7 класс}

\problem{7.1}Расставьте в левой части равенства знаки арифметических операций и скобки вместо звёздочек так, чтобы равенство стало верным для всех x, отличных от нуля.

$* \ x  * \dfrac{1}{x} * x \ \dfrac{1}{x} * x * \dfrac{1}{x} * x * \dfrac{1}{x} * x \ * = (x+1)^2$\\

Ответ: $x:\dfrac{1}{x} + (x \cdot \dfrac{1}{x} + x \cdot \dfrac{1}{x}) \cdot x + \dfrac{1}{x} \cdot x.$

Решение:\\
Заметим, что $(x+1)^2 = x^2 + 2x + 1$
Теперь расставим знаки и скобки:
$x:\dfrac{1}{x}+ (x\cdot\dfrac{1}{x} + x\cdot\dfrac{1}{x})\cdot x + \dfrac{1}{x}\cdot x = x^2 + (1+1)\cdot x + 1 = x^2 + 2\cdot x + 1 = (x+1)^2$

\hspace{3mm}

Критерии:\\
Любая правильная расстановка знаков и скобок: 7 баллов.\\
Неправильная расстановка знаков и скобок: 0 баллов.

\hspace{3mm}

\problem{7.2}Можно ли покрасить некоторые белые поля доски $2022\times2023$ (изначально все поля белые) в чёрный цвет так, чтобы белых полей было на 2023 меньше, чем чёрных?

\hspace{3mm}

Ответ: Нет, нельзя.

Решение:\\
Предположим, что можно. Тогда пусть белых полей X, а чёрных полей $X + 2023$. Следовательно, суммарное количество полей $X + X + 2023 = 2\cdot X + 2023$, 
где $2\cdot X$ - число чётное при любом натуральном X. 
Значит, число $2\cdot X + 2023$ - нечётное, так как является суммой чётного и нечётного. Получается, общее количество полей на доске нечётно. Но всего полей на доске $2022\times2023$ - чётное количество, так как является произведением чётного и нечётного числа. Противоречие. Следовательно, предположение о том, что это можно было сделать, является неверным. Отсюда следует, что нельзя так покрасить. 

\hspace{3mm}

Критерии:\\
Любое правильное решение с верным обоснованием: \textbf{7 баллов.}\\
Показано, что после покраски чёрных и белых суммарно нечетное количество: \textbf{4 балла.}

\newpage

\hspace{3mm}

\noindentПоказано, что всего полей на доске чётное количество: \textbf{2 балла.}\\
Идеи, продвигающие ход решения, но недостаточные для получения верного ответа: \textbf{1 балл.}\\
Неверное решение или идеи, не продвигающие ход решения: \textbf{0 баллов.}\\
Верный или неверный ответ без обоснования: \textbf{0 баллов.}

\hspace{3mm}

\problem{7.3}Пять футбольных команд участвуют в круговом турнире. Могут ли команды по завершению турнира набрать следующие очки: 10, 9, 6, 4, 3? (В круговом футбольном турнире каждая команда играет с каждой по одному матчу. За победу дается 3 очка, за ничью — 1, за поражение — 0).

\hspace{3mm}

Решение:\\
Всего было сыграно $\dfrac{5\cdot4}{2} = 10$ матчей. Пусть команда А играет с командой Б. Рассмотрим всевозможные исходы данного матча:

А выиграла - получила 3 очка, Б проиграла - получила 0 очков $\Rightarrow$ Суммарное количество очков всех команд увеличилось на 3.

Б выиграла - получила 3 очка, А проиграла - получила 0 очков $\Rightarrow$ Суммарное количество очков всех команд увеличилось на 3.

А и Б сыграли вничью и получили по 1 очку $\Rightarrow$ Суммарное количество очков всех команд увеличилось на 2.

Заметим, что суммарное количество очков всех команд максимально может увеличиваться на 3 за один матч. Значит, за 10 матчей максимальное суммарное количество очков всех команд равняется $10\cdot3=30$. 
Но суммарное количество очков, которые набрали команды: 10 + 9 + 6 + 4 + 3 = 32 > 30, чего не может быть.

\hspace{3mm}

Критерии:\\
Любое правильное решение с верным обоснованием: \textbf{7 баллов.}\\
Доказано, что максимальное суммарное количество очков всех команд равняется 30: \textbf{5 баллов.}\\
Идеи, продвигающие ход решения, но недостаточные для получения верного ответа: \textbf{2 балла.}\\
Посчитано, что всего было сыграно 10 матчей: \textbf{1 балл.}\\

\newpage

\hspace{3mm}

\noindentНеверное решение или идеи, не продвигающие ход решения: \textbf{0 баллов.}\\
Верный или неверный ответ без обоснования: \textbf{0 баллов.}

\hspace{3mm}


\problem{7.4}Чингиз записал на доске десятизначное число, которое делится на 6, но его лучший друг Анджа решил над ним подшутить и стёр первую и последнюю цифры этого числа. В результате на доске осталось записано число 20232023. Определите все возможные варианты десятизначного числа, которое мог записать Чингиз до шутки Анджи?

Ответ:\\
1) 1202320230\\
2) 4202320230\\
3) 7202320230\\
4) 2202320232\\
5) 5202320232\\
6) 8202320232\\ 
7) 3202320234\\
8) 6202320234\\
9) 9202320234\\ 
10) 1202320236\\
11) 4202320236\\
12) 7202320236\\ 
13) 2202320238\\
14) 5202320238\\
15) 8202320238

\hspace{3mm}

Решение:\\
Так как число Чингиза делилось на 6, то оно должно делиться на 2 и на 3 одновременно. Следовательно, должны выполняться оба признака делимости на 2 и на 3, а именно: последняя цифра числа должна быть чётной, и сумма цифр числа должна делиться на 3.\\
Переберем все варианты:\\
Пусть последняя цифра была 0, тогда сумма всех цифр равна 14. Ближайшие числа, кратные 3, не меньше 14 - это 15, 18, 21, 24 …\\
Для 15 первая цифра должна быть 1.

\newpage

\hspace{3mm}

\noindentДля 18 первая цифра должна быть 4.\\
Для 21 первая цифра должны быть 7.\\
Для 24 первая цифра должна быть 10, чего не может быть и для вариантов больше также не может быть.\\
Следовательно, могли быть числа:\\
1202320230\\
4202320230\\
7202320230

Пусть последняя цифра была 2, тогда сумма всех цифр равна 16. Ближайшие числа, кратные 3, не меньше 16 - это 18, 21, 24, 27…\\
Для 18 первая цифра должна быть 2.\\
Для 21 первая цифра должна быть 5.\\
Для 24 первая цифра должны быть 8.\\
Для 27 первая цифра должна быть 11, чего не может быть и для вариантов больше также не может быть.\\
Следовательно, могли быть числа:\\
2202320232\\
5202320232\\
8202320232 

Пусть последняя цифра была 4, тогда сумма всех цифр равна 18. Ближайшие числа, кратные 3, не меньше 18 - это 18, 21, 24, 27, 30...\\
Для 18 первая цифра должна быть 0, чего не может быть, так как число не может начинаться с 0.\\
Для 21 первая цифра должна быть 3.\\
Для 24 первая цифра должны быть 6.\\
Для 27 первая цифра должна быть 9.\\
Для 30 первая цифра должна быть 12, чего не может быть и для вариантов больше также не может быть.\\
Следовательно, могли быть числа:\\
3202320234\\
6202320234\\
9202320234\\

\newpage

\hspace{3mm}

\noindentПусть последняя цифра была 6, тогда сумма всех цифр равна 20. Ближайшие числа, кратные 3, не меньше 20 - это 21, 24, 27, 30...\\
Для 21 первая цифра должна быть 1.\\
Для 24 первая цифра должны быть 4.\\
Для 27 первая цифра должна быть 7.\\
Для 30 первая цифра должна быть 10, чего не может быть и для вариантов больше также не может быть.\\
Следовательно, могли быть числа:\\
1202320236\\
4202320236\\
7202320236 

Пусть последняя цифра была 8, тогда сумма всех цифр равна 22. Ближайшие числа, кратные 3, не меньше 22 - это 24, 27, 30, 33...\\
Для 24 первая цифра должна быть 2.\\
Для 27 первая цифра должны быть 5.\\
Для 30 первая цифра должна быть 8.\\
Для 33 первая цифра должна быть 11, чего не может быть и для вариантов больше также не может быть.\\
Следовательно, могли быть числа:\\
2202320238\\
5202320238\\
8202320238 

Выпишем все возможные варианты:\\
1) 1202320230\\
2) 4202320230\\
3) 7202320230\\
4) 2202320232\\
5) 5202320232\\
6) 8202320232\\ 
7) 3202320234\\
8) 6202320234\\
9) 9202320234\\ 
10) 1202320236

\newpage

\hspace{3mm}

\noindent11) 4202320236\\
12) 7202320236\\ 
13) 2202320238\\
14) 5202320238\\
15) 8202320238 
 

Критерии:\\
15 правильных чисел со всеми обоснованиями: \textbf{7 баллов.}\\
10-14 правильных чисел со всеми обоснованиями: \textbf{6 баллов.}\\
15 правильных чисел без обоснований: \textbf{5 баллов.}\\
10-14 правильных чисел без обоснований: \textbf{4 балла.}\\
5-9 правильных чисел со всеми обоснованиями: \textbf{4 балла.}\\
5-9 правильных чисел без обоснований: \textbf{3 балла.}\\
1-4 правильных чисел с обоснованием: \textbf{2 балла.}\\
1-4 правильных чисел без обоснования: \textbf{1 балл.}\\
0 правильных чисел с хорошими идеями, помогающие решению: \textbf{1 балл.}\\
0 правильных чисел с идеями, не помогающие решению или неправильное решение: \textbf{0 баллов.}

\hspace{3mm}

\problem{7.5}В Элисту приехало 2023 жителя острова ВсОШ, и они сразу же встали в огромный круг друг за другом. Первый сказал: «Следующий за мной Второй является лжецом». Второй сказал: «Следующий за мной Третий является лжецом». Третий сказал: «Следующий за мной Четвёртый является лжецом» и так далее до 2023-го, который сказал: «Следующий за мной Первый является лжецом». Известно, что каждый житель острова - либо рыцарь, который всегда говорит правду, либо лжец, который всегда лжёт, либо хитрец, который может говорить правду или лгать по своему усмотрению. Какое минимальное количество хитрецов могло быть среди 2023-х приезжих жителя острова ВсОШ? 

\hspace{3mm}

Ответ: 1 хитрец.

Решение:\\
\textit{Оценка:} Пусть среди приезжих нету хитрецов.

\newpage

\hspace{3mm}

\noindentТогда рассмотрим всевозможные варианты, кем мог быть Первый:\\
1. Предположим, что Первый был рыцарем. Следовательно, Второй был лжецом, а значит, Третий снова был рыцарем, Четвёртый - это лжец и так далее. Каждый нечётный в кругу является рыцарем, а каждый чётный - лжецом. Значит, 2023-ий и Первый оба являются рыцарями, но 2023-ий лжёт, говоря, что Первый является лжецом, но рыцарь не может лгать. Противоречие.\\
2. Предположим, что Первый был лжецом. Следовательно, Второй был рыцарем, а значит, Третий снова был лжецом, Четвёртый - это рыцарь и так далее. Каждый нечётный в кругу является лжецом, а каждый чётный - рыцарем. Значит, 2023-ий и Первый оба являются лжецами, но 2023-ий сказал, что Первый является лжецом, то есть он сказал правду, а лжец не может говорить правду. Противоречие.\\
Следовательно, среди приезжих как минимум один хитрец.

\textit{Пример на одного хитреца}: Первый рыцарь, Второй лжец, Третий рыцарь, Четвёртый лжец … 2021-ый рыцарь, 2022-ой лжец, 2023-ий хитрец.

\hspace{3mm}

Критерии:\\
Приведена и доказана правильная оценка на одного хитреца и приведён любой правильный пример с одним хитрецом: \textbf{7 баллов.}\\
Приведена и доказана правильная оценка на одного хитреца, но пример не приведён или приведён неправильный: \textbf{5 баллов.}\\
Приведён любой правильный пример с одним хитрецом, но не доказано, почему не может быть нуля хитрецов: \textbf{2 балла.}\\
Идеи, продвигающие ход решения, но недостаточные для получения верного ответа: \textbf{2 балла.}\\
Приведена оценка не на одного хитреца или приведён пример не с одним хитрецом: \textbf{0 баллов.}\\
Верный или неверный ответ без обоснования: \textbf{0 баллов.}
















